\documentclass[11pt]{article}

\usepackage[margin=1in]{geometry}
\usepackage{amsmath,amssymb,amsthm}
\usepackage{mathtools}
\usepackage{microtype}
\usepackage{booktabs}
\usepackage{hyperref}
\hypersetup{colorlinks=true,linkcolor=blue,citecolor=blue,urlcolor=blue}

\newcommand{\dynset}{\mathcal{D}}      % dynamic-species index set
\newcommand{\anset}{\mathcal{A}}       % analytic-counterion index set
\DeclareMathOperator{\diag}{diag}
\newcommand{\dd}{\mathrm{d}}

\title{Derivation of the Analytic Bikerman Counterion Closure}
\author{PNPInverse / Forward solver}
\date{\today}

\begin{document}
\maketitle

\begin{abstract}
We derive the analytic concentration profile used in
\texttt{Forward/bv\_solver/boltzmann.py} for inert supporting-electrolyte
counterions. Starting from the electrochemical potential of an ideal
dilute species we obtain the classical Boltzmann distribution
$c = c_b \exp(-z\phi)$. We then replace the ideal entropy with the
Bikerman lattice-gas correction $-\ln\theta$, derive the resulting
algebraic closure for a single inert counterion in equilibrium with a
prescribed potential $\phi(x)$ and a set of dynamic species $\{c_i(x)\}$,
and finally generalize to several inert counterions sharing a common
packing fraction. The result is the steady-state algebraic reduction of
the dynamic Bikerman--Poisson--Nernst--Planck (BPNP) problem in which
the counterion equation is eliminated exactly.
\end{abstract}

%==============================================================
\section{Setup and notation}
%==============================================================

\paragraph{Nondimensionalization.}
All quantities below are nondimensional with the same scales used in the
solver. Potentials are scaled by the thermal voltage
$V_T = k_B T / e$, concentrations by a reference $c_{\mathrm{ref}}$, and
steric coefficients by $1/c_{\mathrm{ref}}$. Dropping tildes, we write
\begin{equation}
\phi \;\leftarrow\; \frac{e\,\phi^{\mathrm{dim}}}{k_B T},
\qquad
c \;\leftarrow\; \frac{c^{\mathrm{dim}}}{c_{\mathrm{ref}}},
\qquad
a \;\leftarrow\; a^{\mathrm{dim}}\, c_{\mathrm{ref}}.
\end{equation}
The Boltzmann factor for an ion of charge $z$ in potential $\phi$ is
$\exp(-z\phi)$ (no explicit $k_BT$). The product $a\,c$ is the local
volume fraction occupied by that species.

\paragraph{Species partition.}
We split the chemical species into
\begin{align}
\dynset &= \{i = 1,\dots,N_d\} \quad
  \text{dynamic species (each has its own NP equation)},\\
\anset &= \{k = 1,\dots,N_a\} \quad
  \text{inert analytic counterions (no electrode reaction, no flux)}.
\end{align}
In the production stack $\dynset = \{\mathrm{O}_2,\mathrm{H}_2\mathrm{O}_2,
\mathrm{H}^+\}$ and $\anset \subseteq \{\mathrm{K}^+,
\mathrm{SO}_4^{2-}, \mathrm{Cs}^+, \mathrm{OH}^-\}$. For each $k\in\anset$
the bulk concentration $c_{b,k}$, charge $z_k$, and steric coefficient
$a_k$ are fixed inputs; the spatially varying $\phi(x)$ is treated as
given by the outer Poisson solve.

\paragraph{Bulk packing fractions.}
Let
\begin{align}
A_{\mathrm{dyn}}(x) &= \sum_{i\in\dynset} a_i\, c_i(x),
   \qquad A^{\mathrm{bulk}}_{\mathrm{dyn}} = \sum_{i\in\dynset} a_i\, c_{b,i}, \\
A_{\mathrm{an}}^{\mathrm{bulk}}  &= \sum_{k\in\anset} a_k\, c_{b,k},
   \qquad
\theta_b = 1 - A^{\mathrm{bulk}}_{\mathrm{dyn}} - A_{\mathrm{an}}^{\mathrm{bulk}}.
\label{eq:thetab}
\end{align}
The bulk free-volume fraction $\theta_b$ is a constant determined entirely
by configuration. The closure below is well posed only if $\theta_b > 0$,
which is enforced at config-parse time.

%==============================================================
\section{Ideal Boltzmann distribution}
%==============================================================

For a dilute species in a potential $\phi$ the electrochemical potential
(in nondimensional form) is~\cite{Newman2004}
\begin{equation}
\mu_i^{\mathrm{id}}(x) \;=\; \ln c_i(x) \;+\; z_i\,\phi(x).
\label{eq:mu_ideal}
\end{equation}
The standard-state contribution $\mu_i^{(0)}$ is suppressed because it
drops out of every equation we use.

\paragraph{Equilibrium of an inert species.}
An ``inert'' species $k$ has
\begin{enumerate}
\item no source or sink anywhere in the domain,
\item zero normal flux on every non-bulk boundary,
\item Dirichlet $c_k = c_{b,k}$ on the bulk boundary, where $\phi = 0$.
\end{enumerate}
Its Nernst--Planck flux is $\mathbf{J}_k = -D_k\,c_k\,\nabla \mu_k$.
At steady state the divergence vanishes and the boundary conditions
above force $\mathbf{J}_k \equiv \mathbf{0}$. With $c_k > 0$ this gives
\begin{equation}
\nabla \mu_k(x) = 0
\qquad\Longrightarrow\qquad
\mu_k(x) = \mu_k^{\mathrm{bulk}}.
\label{eq:zero_grad}
\end{equation}

\paragraph{Boltzmann profile.}
Substituting \eqref{eq:mu_ideal} into \eqref{eq:zero_grad},
\begin{equation}
\ln c_k(x) + z_k\,\phi(x) \;=\; \ln c_{b,k} + z_k\cdot 0,
\end{equation}
so
\begin{equation}
\boxed{\;c_k(x) \;=\; c_{b,k}\,\exp\!\bigl(-z_k\,\phi(x)\bigr).\;}
\label{eq:boltzmann_ideal}
\end{equation}
This is the legacy \texttt{steric\_mode="ideal"} closure. It is unbounded:
for an anion ($z_k=-1$) at large anodic $\phi$, $c_k$ grows
exponentially, eventually exceeding any physical packing limit.

%==============================================================
\section{Bikerman chemical potential}
%==============================================================

The ideal entropy $\ln c_i$ in \eqref{eq:mu_ideal} assumes ions are point
particles. Bikerman's lattice-gas correction accounts for the finite
volume each ion occupies by treating the solvent + ions as a fully
packed lattice with site fraction $a_i$ per ion. Minimizing the
resulting free energy~\cite{Tresset2008} gives the modified chemical
potential
\begin{equation}
\mu_i(x) \;=\; \ln c_i(x) \;+\; z_i\,\phi(x) \;-\; \ln\theta(x),
\label{eq:mu_bikerman}
\end{equation}
where
\begin{equation}
\theta(x) \;=\; 1 - \sum_{j\in\dynset\cup\anset} a_j\,c_j(x)
\label{eq:theta_def}
\end{equation}
is the local free-volume fraction. The extra term $-\ln\theta$ is the
excluded-volume contribution to the free energy: as ions pack into a
region $\theta \to 0$, $\mu_i\to +\infty$, and the species is
electrostatically expelled.

\paragraph{Sign convention.}
The codebase uses the \emph{sign-corrected} form
$\mu_{\mathrm{steric}} = -\ln\theta$ consistently in
\texttt{forms\_logc.py} (\texttt{forms\_logc.py:402}). The
counterion closure below assumes the same sign; mixing the legacy
$+\ln\theta$ branch (still present in \texttt{Forward/dirichlet\_solver.py}
and \texttt{robin\_solver.py}, neither of which is in the production
dispatch) with the closure here would be inconsistent.

%==============================================================
\section{Single-counterion algebraic closure}
%==============================================================

Consider one inert counterion ($N_a = 1$, label it ``$-$''). Equation
\eqref{eq:zero_grad} together with \eqref{eq:mu_bikerman} gives
\begin{equation}
\ln c_-(x) + z_-\,\phi(x) - \ln\theta(x)
\;=\; \ln c_{b,-} - \ln\theta_b,
\label{eq:mu_match}
\end{equation}
i.e.
\begin{equation}
\frac{c_-(x)\,\exp\!\bigl(z_-\,\phi(x)\bigr)}{\theta(x)}
\;=\; \frac{c_{b,-}}{\theta_b}.
\label{eq:zero_flux_alg}
\end{equation}
Introduce the convenient abbreviations
\begin{equation}
q(x) \;=\; \exp\!\bigl(-z_-\,\phi(x)\bigr),
\qquad
B(x) \;=\; \frac{c_{b,-}\,q(x)}{\theta_b}.
\end{equation}
Then \eqref{eq:zero_flux_alg} reads $c_- = B\,\theta$. Splitting
$\theta = 1 - A_{\mathrm{dyn}} - a_-\,c_-$ and isolating $c_-$:
\begin{align}
c_- &= B\bigl(1 - A_{\mathrm{dyn}} - a_- c_-\bigr), \\
c_-\,(1 + a_-\,B) &= B\,(1 - A_{\mathrm{dyn}}), \\
\boxed{\;
c_-(x) \;=\; \frac{c_{b,-}\,q(x)\,\bigl(1-A_{\mathrm{dyn}}(x)\bigr)}
                  {\theta_b + a_-\,c_{b,-}\,q(x)}.
\;}
\label{eq:single_closure}
\end{align}
This is the formula implemented in
\texttt{build\_steric\_boltzmann\_expressions}
(\texttt{Forward/bv\_solver/boltzmann.py:254}).

\paragraph{Limits.}
Three checks recover the expected physics.
\begin{enumerate}
\item \emph{Bulk recovery.} At $\phi = 0$, $q = 1$, and
$A_{\mathrm{dyn}} = A^{\mathrm{bulk}}_{\mathrm{dyn}}$, so
\begin{equation}
c_-^{\mathrm{bulk}}
  = \frac{c_{b,-}\,(1 - A^{\mathrm{bulk}}_{\mathrm{dyn}})}
         {\theta_b + a_-\,c_{b,-}}
  = c_{b,-},
\end{equation}
since $\theta_b + a_-\,c_{b,-} = 1 - A^{\mathrm{bulk}}_{\mathrm{dyn}}$
(\eqref{eq:thetab} with $N_a=1$).

\item \emph{Dilute limit.} If $a_- \to 0$ and $A_{\mathrm{dyn}} \to 0$
then $\theta_b \to 1$ and \eqref{eq:single_closure} collapses to
$c_- = c_{b,-}\,q = c_{b,-}\exp(-z_-\phi)$, the ideal Boltzmann profile
\eqref{eq:boltzmann_ideal}.

\item \emph{High-field saturation.} For an anion at large anodic
$\phi$, $q\to\infty$ and the denominator is dominated by
$a_-\,c_{b,-}\,q$, giving
\begin{equation}
c_-(x) \;\longrightarrow\;
\frac{1 - A_{\mathrm{dyn}}(x)}{a_-}.
\label{eq:saturation}
\end{equation}
The counterion fills exactly the free volume left by the dynamic
species, never more. Contrast with the unbounded ideal Boltzmann, where
$c_- \sim c_{b,-}\,q$ diverges.
\end{enumerate}

\paragraph{Total packing.}
Substituting \eqref{eq:single_closure} into
$\theta = 1 - A_{\mathrm{dyn}} - a_- c_-$ yields
\begin{equation}
\theta(x) \;=\;
\bigl(1 - A_{\mathrm{dyn}}(x)\bigr)\,
\frac{\theta_b}{\theta_b + a_-\,c_{b,-}\,q(x)},
\label{eq:theta_single}
\end{equation}
which is strictly positive whenever $1-A_{\mathrm{dyn}} > 0$ and
$\theta_b > 0$.

%==============================================================
\section{\texorpdfstring{Multi-counterion shared-$\theta$ closure}{Multi-counterion shared-theta closure}}
\label{sec:multi}
%==============================================================

When $|\anset| > 1$ each analytic counterion satisfies its own
zero-flux equilibrium, but all of them share the same packing
fraction $\theta$. Writing
\eqref{eq:zero_flux_alg} for each $k\in\anset$ and abbreviating
\begin{equation}
q_k(x) = \exp(-z_k\,\phi(x)),
\qquad
B_k(x) = \frac{c_{b,k}\,q_k(x)}{\theta_b},
\end{equation}
we have one algebraic equation per counterion,
\begin{equation}
c_k \;=\; B_k\,\theta
\;=\; B_k\Bigl(1 - A_{\mathrm{dyn}} - \sum_{k'\in\anset} a_{k'}\,c_{k'}\Bigr).
\label{eq:multi_system}
\end{equation}

\paragraph{Reduction via the analytic packing sum.}
Let $S(x) = \sum_{k\in\anset} a_k\,c_k(x)$. Multiplying
\eqref{eq:multi_system} by $a_k$ and summing,
\begin{equation}
S = \Bigl(\sum_k a_k B_k\Bigr)\,
   (1 - A_{\mathrm{dyn}} - S),
\end{equation}
hence
\begin{equation}
S \;=\; (1 - A_{\mathrm{dyn}})\,
\frac{\sum_k a_k B_k}{1 + \sum_k a_k B_k}.
\end{equation}
The remaining free volume is
\begin{equation}
1 - A_{\mathrm{dyn}} - S
\;=\; \frac{1 - A_{\mathrm{dyn}}}{1 + \sum_k a_k B_k}
\;=\; \frac{\theta_b\,(1-A_{\mathrm{dyn}})}
           {\theta_b + \sum_k a_k\,c_{b,k}\,q_k(x)}.
\label{eq:shared_freevol}
\end{equation}
Substituting back into \eqref{eq:multi_system} gives the closed-form
expression for each counterion:
\begin{equation}
\boxed{\;
c_k(x) \;=\;
\frac{c_{b,k}\,q_k(x)\,\bigl(1-A_{\mathrm{dyn}}(x)\bigr)}
     {\theta_b + \sum_{k'\in\anset} a_{k'}\,c_{b,k'}\,q_{k'}(x)},
\qquad k\in\anset.
\;}
\label{eq:multi_closure}
\end{equation}
With $A_{\mathrm{dyn}}\equiv 0$ this collapses to the standard
full-equilibrium asymmetric-Bikerman ion-distribution formula
of~\cite{Tresset2008}; the $(1-A_{\mathrm{dyn}})$ pre-factor is the
contribution of the dynamic species' local packing, retained because
they obey their own NP equations rather than being in equilibrium
with $\phi$. The relation is worked out in detail in
\S\ref{sec:tresset}.

\paragraph{Structure of the closure.}
\begin{itemize}
\item The numerator $c_{b,k}\,q_k\,(1-A_{\mathrm{dyn}})$ is
\emph{per-ion}: each counterion has its own Boltzmann factor and the
common dynamic-species free-volume multiplier.
\item The denominator
$\theta_b + \sum_{k'} a_{k'}\,c_{b,k'}\,q_{k'}$ is \emph{shared} across
all $k$. This is exactly what makes the multi-ion problem decouple:
the algebra never requires solving a coupled local nonlinear system.
\item For $|\anset| = 1$ the denominator reduces to
$\theta_b + a_-\,c_{b,-}\,q$ and \eqref{eq:multi_closure} collapses to
\eqref{eq:single_closure}.
\item For two oppositely charged ions (e.g.\ $\mathrm{K}^+$ with
$z=+1$ and $\mathrm{SO}_4^{2-}$ with $z=-2$) the denominator contains
both $\exp(-\phi)$ and $\exp(+2\phi)$, so it is bounded away from zero
on both sides of $\phi = 0$.
\end{itemize}

\paragraph{High-field saturation, multi-ion.}
For large $\phi$ of either sign, one term in
$\sum_{k'} a_{k'}\,c_{b,k'}\,q_{k'}$ dominates. Call its index
$k^\star = \arg\max_{k'} (-z_{k'}\,\phi)$. Then
\begin{align}
c_{k^\star}(x) &\;\longrightarrow\;
   \frac{1 - A_{\mathrm{dyn}}(x)}{a_{k^\star}}, \\
c_k(x) &\;\longrightarrow\; 0
\quad \text{for } k\neq k^\star,
\end{align}
i.e. the most strongly attracted analytic counterion fills all of the
remaining free volume and the others are depleted. This is the multi-ion
analog of \eqref{eq:saturation}.

%==============================================================
\section{Relation to Tresset's generalized Poisson--Fermi formula}
\label{sec:tresset}
%==============================================================

The full-equilibrium limit ($A_{\mathrm{dyn}} \equiv 0$) of
\eqref{eq:multi_closure} is not new: it coincides with Tresset's
generalized Fermi-like ion distribution for size-asymmetric
multi-species lattice gases~\cite[Eq.~(19), page 061506-4]{Tresset2008}.
The hybrid extension we use here --- closing only the inert subset
analytically while the reactive species remain dynamic --- amounts to
one additional algebraic step that introduces the explicit
$(1 - A_{\mathrm{dyn}}(x))$ pre-factor.

\paragraph{Tresset's full-equilibrium formula.}
Minimizing the lattice-gas free energy for $M$ species of
valence $z_i$ and excluded volume $v_i$ and imposing thermodynamic
equilibrium $\mu_i(\mathbf{r}) = \mu_i^{\mathrm{bulk}}$ for
\emph{every} species, Tresset~\cite[Eq.~(19)]{Tresset2008} obtains
\begin{equation}
\phi_i(x)
\;=\;
\frac{\exp\!\bigl(-z_i\,\psi^*(x)\bigr)\,\phi_i^\infty}
     {\phi_{\mathrm{FS}}^\infty +
      \sum_{j=1}^M \exp\!\bigl(-z_j\,\psi^*(x)\bigr)\,\phi_j^\infty},
\qquad 1 \le i \le M,
\label{eq:tresset19}
\end{equation}
where $\phi_i = c_i\,v_i$ is the local volume fraction of species
$i$, $\phi_i^\infty = c_i^\infty v_i$ its bulk value,
$\phi_{\mathrm{FS}}^\infty = 1 - \sum_j \phi_j^\infty$ the bulk
free-space volume fraction, and $\psi^* = e\psi/k_BT$ the
dimensionless electrostatic potential.

\paragraph{Notation map.}
\begin{center}
\begin{tabular}{ll}
\toprule
Tresset~\cite{Tresset2008} & This work \\
\midrule
$v_i$ (excluded volume) & $a_i$ (nondim steric coefficient) \\
$\phi_i = c_i\,v_i$ & $a_i\,c_i$ \\
$\phi_i^\infty = c_i^\infty\,v_i$ & $a_i\,c_{b,i}$ \\
$\phi_{\mathrm{FS}}^\infty = 1 - \sum_j \phi_j^\infty$ & $\theta_b$ \eqref{eq:thetab} \\
$\psi^*(x) = e\psi(x)/k_BT$ & $\phi(x)$ (already nondim) \\
$\exp(-z_i\,\psi^*)$ & $q_i = \exp(-z_i\,\phi)$ \\
\bottomrule
\end{tabular}
\end{center}
\emph{Notation caveat:} Tresset uses $\phi$ for volume fraction and
$\psi$ for potential; we use $\phi$ for potential. The collision is
unfortunate but the dictionary above is unambiguous.

\paragraph{Equivalence on the full-equilibrium subset.}
Set $A_{\mathrm{dyn}}\equiv 0$ in \eqref{eq:multi_closure}, multiply
both sides by $a_k$, and apply the dictionary:
\begin{equation}
a_k\,c_k(x)
\;=\;
\frac{a_k\,c_{b,k}\,q_k(x)}
     {\theta_b + \sum_{k'} a_{k'}\,c_{b,k'}\,q_{k'}(x)}
\;\;\overset{\text{map}}{\longleftrightarrow}\;\;
\phi_k(x)
\;=\;
\frac{\phi_k^\infty\,e^{-z_k\psi^*}}
     {\phi_{\mathrm{FS}}^\infty + \sum_{k'} \phi_{k'}^\infty\,e^{-z_{k'}\psi^*}},
\end{equation}
which is \eqref{eq:tresset19} verbatim. The algebraic steps in
Tresset (his Eqs.~(15)--(19), pages 061506-3 to 061506-4) coincide
with \S\ref{sec:multi} of this note: identical chemical potential
$\mu_i = \ln c_i + z_i\phi - \ln\theta$, identical equilibrium
condition $\nabla\mu_i = 0$, identical algebraic inversion.

\paragraph{Modification for hybrid coupling.}
Tresset's \eqref{eq:tresset19} assumes \emph{every} species is in
equilibrium with $\phi$. In the modified-PNP problem the dynamic
species $i \in \dynset$ are not: they obey Nernst--Planck equations
with electrode-reaction source terms, and their profiles
$c_i(x)$ are determined by transport and kinetics. The zero-flux
argument $\nabla\mu_i = 0$ \eqref{eq:zero_grad} \emph{fails for
them}; it still holds for the inert subset $k \in \anset$.

Repeating Tresset's algebra on the inert subset only, with the
dynamic species' packing $A_{\mathrm{dyn}}(x) =
\sum_{i\in\dynset} a_i c_i(x)$ treated as a known spatial
background, the equilibrium condition for $k \in \anset$ becomes
\begin{equation}
\ln c_k(x) + z_k\,\phi(x)
\;-\; \ln\bigl[1 - A_{\mathrm{dyn}}(x) - S(x)\bigr]
\;=\; \ln c_{b,k} - \ln\theta_b,
\label{eq:hybrid_equilibrium}
\end{equation}
with $S(x) = \sum_{k\in\anset} a_k\,c_k(x)$ the analytic-counterion
packing. The only difference from Tresset's intermediate step is
that the local free-space fraction
$\phi_{\mathrm{FS}}(x) = 1 - \sum_j \phi_j(x)$ has been split into
a dynamic part $A_{\mathrm{dyn}}(x)$ (now a spatial background, not
a sum-of-equilibrium-profiles) and an inert part $S(x)$ (the
unknown). Following exactly the same multiply-by-$a_k$/sum/isolate
procedure as in \S\ref{sec:multi} returns
\begin{equation}
c_k(x) \;=\;
\frac{c_{b,k}\,q_k(x)\,\bigl(1 - A_{\mathrm{dyn}}(x)\bigr)}
     {\theta_b + \sum_{k'\in\anset} a_{k'}\,c_{b,k'}\,q_{k'}(x)},
\qquad k\in\anset,
\label{eq:hybrid_closure}
\end{equation}
which is \eqref{eq:multi_closure}. The lone structural difference
from Tresset's \eqref{eq:tresset19} is the explicit
$\bigl(1 - A_{\mathrm{dyn}}(x)\bigr)$ factor in the numerator.

\paragraph{When the modification matters.}
If the dynamic species' local packing equals its bulk value
everywhere ($A_{\mathrm{dyn}}(x) \equiv A_{\mathrm{dyn}}^{\mathrm{bulk}}$),
the pre-factor is a global constant. Defining $\theta_b' =
\theta_b/(1 - A_{\mathrm{dyn}}^{\mathrm{bulk}})$ and
$c_{b,k}' = c_{b,k}/(1 - A_{\mathrm{dyn}}^{\mathrm{bulk}})$ absorbs
it into rescaled bulk quantities and \eqref{eq:hybrid_closure}
reduces to Tresset's \eqref{eq:tresset19} exactly. The correction
matters precisely when $A_{\mathrm{dyn}}(x)$ varies in space, which
is the rule, not the exception, in PNP problems with electrode
reactions. In the production stack the proton $\mathrm{H}^+$
Boltzmann-piles up at the OHP under cathodic polarization, so
$A_{\mathrm{dyn}}(\mathrm{OHP}) \gg A_{\mathrm{dyn}}^{\mathrm{bulk}}$
and the inert counterion is volumetrically squeezed out of the
diffuse layer.

\paragraph{What is and is not novel.}
Equation \eqref{eq:tresset19} --- the size-asymmetric,
all-equilibrium lattice-gas distribution --- is
Tresset~\cite{Tresset2008}. The hybrid-coupling
extension \eqref{eq:hybrid_closure} = \eqref{eq:multi_closure} ---
closing the inert subset analytically while keeping the reactive
species dynamic, with the explicit $(1 - A_{\mathrm{dyn}}(x))$
pre-factor --- is, to our knowledge, not previously written in
closed form in the modified-PNP literature. The algebraic
modification is one step; the use case (reactive analyte species
coupled to an inert supporting electrolyte through a closed-form
analytic counterion in a steric-aware modified-PNP problem)
appears to be original to this work.

%==============================================================
\section{Poisson coupling and remaining BVP}
%==============================================================

The closure \eqref{eq:multi_closure} eliminates the analytic
counterion(s) from the dynamic problem. What remains is

\paragraph{(i) Poisson with closed-form counterion charge.}
\begin{equation}
-\nabla\!\cdot\!\bigl(\varepsilon\,\nabla\phi\bigr)
\;=\; \rho_{\mathrm{chg}}\Bigl(
        \sum_{i\in\dynset} z_i\,c_i
        \;+\; \sum_{k\in\anset} z_k\,c_k[\phi,\{c_i\}]
      \Bigr),
\end{equation}
where $\rho_{\mathrm{chg}}$ is the nondimensional charge prefactor
($\texttt{charge\_rhs\_prefactor}$) and the $c_k[\phi,\{c_i\}]$ are the
explicit functionals \eqref{eq:multi_closure}.

\paragraph{(ii) Nernst--Planck for dynamic species with full packing.}
For each $i \in \dynset$, the modified Poisson--Nernst--Planck
flux~\cite{Kilic2007b} reads
\begin{align}
\partial_t c_i + \nabla\!\cdot\!\mathbf{J}_i &= R_i,\\
\mathbf{J}_i &= -D_i\, c_i\, \nabla \mu_i,
\qquad
\mu_i = \ln c_i + z_i\phi - \ln\theta,
\end{align}
with the \emph{total} packing
\begin{equation}
\theta(x) \;=\; 1 - A_{\mathrm{dyn}}(x) - \sum_{k\in\anset} a_k\,c_k[\phi,\{c_i\}],
\end{equation}
which by \eqref{eq:shared_freevol} simplifies to
\begin{equation}
\theta(x) \;=\;
\frac{\theta_b\,\bigl(1-A_{\mathrm{dyn}}(x)\bigr)}
     {\theta_b + \sum_{k\in\anset} a_k\,c_{b,k}\,q_k(x)}.
\end{equation}
This is the key cross-check on consistency: \emph{the same} closure
expression must appear in (a)~the Poisson charge density and
(b)~the dynamic-species steric chemical potential. Wiring them from
distinct expressions would amount to solving two different models, and
was the failure mode that prompted the refactor of the original
ideal-Boltzmann path (see
\texttt{docs/solver/steric\_analytic\_clo4\_reduction\_handoff.md}).

%==============================================================
\section{Equivalence to dynamic Bikerman--NP}
%==============================================================

The reduction \eqref{eq:multi_closure} is \emph{not} an approximation:
it is the exact steady-state solution of the dynamic Bikerman--NP
equation for each inert counterion, provided
\begin{enumerate}
\item the system is at steady state,
\item each counterion has no electrode reaction, no source/sink, and
no imposed flux on any non-bulk boundary,
\item the bulk boundary fixes $c_k = c_{b,k}$ and $\phi = 0$,
\item the same Bikerman chemical potential
\eqref{eq:mu_bikerman} acts on dynamic and counterion species, and
\item the total packing $\theta$ in the dynamic species' residual
includes the analytic counterion's contribution.
\end{enumerate}
Under these assumptions an explicit calculation shows
$\mathbf{J}_k = 0$ for the closure profile, so it is the unique
solution of the counterion's Nernst--Planck equation up to a
multiplicative gauge that the bulk Dirichlet condition fixes.

The reduction is \emph{not} equivalent in any of the following
situations and should not be used: (a) transient double-layer charging
with finite analyte flux through the supporting electrolyte;
(b)~ion-specific adsorption or chemisorption on the electrode;
(c)~models with a finite closest-approach plane treated dynamically
rather than as a Stern boundary condition.

%==============================================================
\section{Numerical implementation notes}
%==============================================================

Three numerical guards in the implementation are worth noting because
they appear in the production UFL but are not part of the formal
derivation.

\paragraph{Potential clamp.} The Boltzmann factor $q_k = \exp(-z_k\phi)$
is replaced with $\exp(-z_k\,\widetilde\phi)$ where
$\widetilde\phi = \min(\max(\phi,-\phi_{\mathrm{clamp},k}),
\phi_{\mathrm{clamp},k})$. This is a Newton-stability guard during
continuation; on a converged solution $\widetilde\phi \equiv \phi$
provided $|\phi|$ stays inside the clamp. The clamp is per-ion to allow
asymmetric bounds. The shared denominator
in \eqref{eq:multi_closure} uses each ion's own clamp for its own term,
which is consistent because each $\phi_{\mathrm{clamp},k}$ only ever
appears multiplied by $z_k$ inside the $k$-th summand.

\paragraph{Free-volume floor.} The factor $1-A_{\mathrm{dyn}}$ in the
numerator is replaced by $\max(1-A_{\mathrm{dyn}},\,\epsilon_{\mathrm{floor}})$
with $\epsilon_{\mathrm{floor}} = 10^{-10}$. This is a finite-element
interpolation safety: even with $\theta > 0$ pointwise in the analytic
solution, the FE basis can produce small negative $(1-A_{\mathrm{dyn}})$
values on coarse meshes inside Newton iterates. The floor is never
active on a converged, physically meaningful solution.

\paragraph{$z$-ramp scaling.} The Poisson charge contribution
$z_k\,c_k[\phi,\{c_i\}]$ is multiplied by a shared scaling Function
$z_{\mathrm{scale}} \in [0,1]$. Continuation strategies that ramp
dynamic-species charges from $0$ to $z_i$ also ramp the analytic
counterion's charge through this factor, keeping the bulk
electroneutral throughout the homotopy. Default $z_{\mathrm{scale}} = 1$
recovers the physical model.

%==============================================================
\section{Code map}
%==============================================================

\begin{center}
\begin{tabular}{ll}
\toprule
Mathematical object & Location \\
\midrule
$\mu_i = \ln c_i + z_i\phi - \ln\theta$
   & \texttt{Forward/bv\_solver/forms\_logc.py:402} \\
$c_k$ \eqref{eq:multi_closure}
   & \texttt{boltzmann.py:254} (\texttt{c\_steric\_expr}) \\
$a_k\,c_k$ packing contribution
   & \texttt{boltzmann.py:257} (\texttt{packing\_contribution}) \\
$z_k\,c_k$ Poisson source
   & \texttt{boltzmann.py:258} (\texttt{charge\_density}) \\
$\theta_b > 0$ validation
   & \texttt{boltzmann.py:159--171} \\
Double-counting validation
   & \texttt{boltzmann.py:176--201} \\
Ideal closure $c_k = c_{b,k}\exp(-z_k\phi)$
   & \texttt{boltzmann.py:380--382} \\
\bottomrule
\end{tabular}
\end{center}

\paragraph{Configuration entry.}
Each analytic counterion is declared as a dict in
\texttt{bv\_bc.boltzmann\_counterions}, e.g.\
\begin{verbatim}
{
  "z": +1,
  "c_bulk_nondim": C_KPLUS_HAT,
  "a_nondim":      A_KPLUS_HAT,
  "phi_clamp":     50.0,
  "steric_mode":   "bikerman",
}
\end{verbatim}
with \texttt{steric\_mode="ideal"} selecting the legacy closure
\eqref{eq:boltzmann_ideal} and \texttt{steric\_mode="bikerman"}
selecting \eqref{eq:multi_closure}.

%==============================================================
\begin{thebibliography}{9}

\bibitem{Kilic2007b}
M.~S.~Kilic, M.~Z.~Bazant, and A.~Ajdari.
\newblock Steric effects in the dynamics of electrolytes at large applied
voltages.\ {II}.~{M}odified {P}oisson--{N}ernst--{P}lanck equations.
\newblock \emph{Physical Review E}, 75(2):021503, 2007.
\newblock \url{https://doi.org/10.1103/PhysRevE.75.021503}
\newblock arXiv preprint: \url{https://arxiv.org/abs/physics/0611232}

\bibitem{Tresset2008}
G.~Tresset.
\newblock Generalized {P}oisson--{F}ermi formalism for investigating size
correlation effects with multiple ions.
\newblock \emph{Physical Review E}, 78(6):061506, 2008.
\newblock \url{https://doi.org/10.1103/PhysRevE.78.061506}

\bibitem{Newman2004}
J.~Newman and K.~E.~Thomas-Alyea.
\newblock \emph{Electrochemical Systems}.
\newblock Wiley-Interscience, Hoboken NJ, 3rd edition, 2004.

\end{thebibliography}

\end{document}
